% ═══════════════════════════════════════════════════════════════════════════
% Chapter 1: Introduction
% ═══════════════════════════════════════════════════════════════════════════

\chapter{Introduction} \label{chap:introduction}

\epigraph{The most practical thing is the sound theory.}{\textmd{Werner Heisenberg}}

\section{Motivation}

Velocity-Based Training (VBT) has emerged as a scientifically rigorous methodology for strength training prescription in the past decade. Unlike traditional percentage-based training that prescribes loads as fixed percentages of one-repetition maximum (1RM), VBT uses the instantaneous velocity of the barbell as the primary variable for training prescription and monitoring. This approach offers several critical advantages:

\begin{enumerate}
    \item \textbf{Daily Auto-regulation:} An athlete's 1RM fluctuates daily based on fatigue, hydration, sleep quality, and psychological state. Training at 75\% of a theoretically determined 1RM may represent a stimulus ranging from intensely fatiguing to insufficiently challenging depending on the athlete's readiness on that particular day. VBT addresses this by using velocity as an objective proxy for relative intensity.

    \item \textbf{Velocity-Fatigue Relationship:} Research has established that velocity loss within a set correlates strongly with neuromuscular fatigue. Terminating sets at specific velocity loss thresholds (e.g., 20\% loss from first rep) produces superior adaptations compared to training to muscular failure while reducing injury risk and central nervous system fatigue.

    \item \textbf{Exercise-Specific Velocity Zones:} Different exercises exhibit distinct velocity profiles for equivalent relative intensities. A back squat at 85\% 1RM will have a different mean velocity than a bench press at the same relative load. Direct velocity measurement enables exercise-specific zone targeting.
\end{enumerate}

The accuracy requirements for VBT devices are stringent. A deviation of even 0.05 m/s in mean concentric velocity (MCV) measurement can shift the training stimulus across distinct physiological adaptation zones, fundamentally altering the intended training effect. This thesis addresses the complete engineering challenge of building a measurement system that meets these requirements while providing the transparency and reproducibility necessary for academic research.

\section{Problem Statement}

The commercial VBT measurement ecosystem consists of three primary modalities, each with significant limitations for research applications:

\subsection{Linear Position Transducers (LPTs)}

Linear Position Transducers (e.g., GymAware, Tendo Unit) measure cable displacement as the barbell moves. While historically considered the ``gold standard'' due to their high accuracy (typically $\pm$0.5 cm), LPTs suffer from fundamental limitations:

\begin{itemize}
    \item \textbf{Barpath Constraint:} The physical tether constrains the barbell to predominantly vertical motion, altering the natural bar path during explosive multi-joint lifts such as the snatch and clean where significant horizontal displacement is biomechanically essential.

    \item \textbf{One-Dimensional Measurement:} LPTs provide only vertical displacement, failing to capture barbell tilt, lateral drift, or rotational dynamics that are clinically relevant for technique analysis.

    \item \textbf{Practical Limitations:} The tether creates setup complexity, can become a tripping hazard, and limits the range of exercises that can be measured.
\end{itemize}

\subsection{Camera-Based Systems}

Marker-based optoelectronic motion capture systems (Vicon, Qualisys, OptiTrack) provide sub-millimeter accuracy and represent the true criterion standard for biomechanical research. However:

\begin{itemize}
    \item \textbf{Cost:} Multi-camera systems typically cost \$50,000-\$250,000, placing them outside reach of most training facilities and individual researchers.

    \item \textbf{Laboratory Footprint:} Fixed installation requires dedicated calibrated volumes, precluding field deployment in gyms or outdoor settings.

    \item \textbf{Calibration Complexity:} Extensive setup time (>30 minutes for multi-camera calibration) and specialized operator knowledge required.
\end{itemize}

Video-based systems using smartphone feeds (e.g., ``My Lift'' application) offer accessibility but suffer from rolling-shutter distortion, limited frame rates (30-60 fps), and sensitivity to lighting conditions.

\subsection{Wearable IMU Systems}

Commercial IMU-based VBT devices (PUSH Band, Vmaxpro, Beast Sensor) solve portability problems but introduce severe limitations:

\begin{itemize}
    \item \textbf{Black Box Algorithms:} Proprietary sensor-fusion algorithms with undocumented filtering, undocumented sampling rates, and unquantified latency make research validation impossible.

    \item \textbf{Filtered Outputs:} Most devices stream only velocity summaries, providing no access to raw inertial data required for algorithm development and validation.

    \item \textbf{Accuracy Concerns:} Literature reports Limits of Agreement as wide as $\pm$0.11 m/s for commercial devices against optical ground truth, an error envelope spanning nearly an entire velocity zone.
\end{itemize}

\subsection{The Fundamental Challenge}

The core challenge of IMU-based VBT stems from the double integration required to derive position from acceleration. Velocity is obtained by integrating linear acceleration; position is obtained by integrating velocity. Any uncompensated bias $b_a$ in the accelerometer produces position error growing \textit{quadratically} with time:

\begin{equation}
    \epsilon_p(t) = \frac{1}{2} b_a \cdot t^2
\end{equation}

For a typical MEMS accelerometer bias of 5 mg (0.049 m/s\textsuperscript{2}), this yields position error of 1.23 meters after only 10 seconds---catastrophic failure for a barbell that may have moved only 0.5 meters per repetition.

\section{Research Objectives}

This dissertation addresses the complete design space of an IMU-based VBT measurement system, from embedded sensor acquisition through real-time signal processing to hardware acceleration architectures. The specific objectives are:

\begin{enumerate}
    \item \textbf{High-Fidelity Data Acquisition:} Capture raw, unfiltered 1 kHz 6-DOF IMU data at $\pm$16 g accelerometer and $\pm$2000 deg/s gyroscope full-scale ranges, preserving high-frequency shock transients of weightlifting movements.

    \item \textbf{Hardware Synchronization:} Achieve sub-millisecond synchronization between IMU and camera streams through dedicated level-shifting circuit and ICM-42688-P's internal FSYNC tagging engine.

    \item \textbf{Wireless Telemetry:} Transmit high-bandwidth IMU data stream from barbell to host PC via ESP-NOW protocol, avoiding physical tethering.

    \item \textbf{Open-Source Processing Pipeline:} Implement mathematically rigorous pipeline with attitude estimation (VQF, ESKF, Madgwick), ZUPT detection, and per-rep batch smoothing achieving sub-5 cm position accuracy IMU-only.

    \item \textbf{Comprehensive Validation:} Validate against optical ground truth across 55 sessions, 622 annotated repetitions, six compound exercises, quantifying accuracy, latency, and reliability.

    \item \textbf{Hardware Acceleration Architecture:} Design FPGA and ASIC implementations achieving 10$\times$+ energy efficiency improvement while maintaining algorithmic fidelity, enabling production deployment.

    \item \textbf{BIDS-Compatible Dataset:} Persist sessions in schema-versioned JSON format with 120+ metadata fields covering anthropometrics, training context, equipment specifications, and subjective assessments.
\end{enumerate}

\section{Thesis Contributions}

The principal contributions of this dissertation are:

\begin{enumerate}
    \item A complete, open-source IMU-based VBT measurement system with documented hardware schematics, firmware source code, and host-side processing pipeline.

    \item A causal, real-time rep-segmentation algorithm achieving sub-50 ms latency per repetition with sub-5 cm position accuracy without camera fusion.

    \item Rigorous validation dataset comprising 55 sessions across six exercises with 622 manually verified repetitions, publicly available for reproducibility.

    \item Complete FPGA/ASIC migration path with RTL specifications, fixed-point analysis, and power/performance estimation for 28nm CMOS implementation.
\end{enumerate}

\section{Document Organization}

This dissertation is organized as follows:

\begin{description}
    \item[Chapter 2:] Background and Literature Review -- VBT principles, sensor technologies, attitude estimation algorithms, and ZUPT strategies.

    \item[Chapter 3:] System Architecture -- High-level design decisions, component selection rationale, and integration strategy.

    \item[Chapter 4:] Hardware Design -- IMU configuration, synchronization circuitry, wireless radio design, PCB layout considerations.

    \item[Chapter 5:] Embedded Firmware -- ESP32 interrupt architecture, SPI driver implementation, ESP-NOW protocol integration, power management.

    \item[Chapter 6:] Host Software Architecture -- C++17 application design, real-time data ingestion, session management, GUI implementation.

    \item[Chapter 7:] Signal Processing Pipeline -- Mathematical foundations of orientation filtering, velocity integration, position estimation, and drift correction.

    \item[Chapter 8:] Rep Segmentation Algorithm -- Streaming boundary detection, back-dating strategy, per-rep batch smoothing, metrics computation.

    \item[Chapter 9:] Validation and Results -- Dataset description, accuracy metrics, latency measurements, comparison against ground truth.

    \item[Chapter 10:] FPGA Implementation -- Architecture mapping, timing analysis, resource utilization, fixed-point quantization.

    \item[Chapter 11:] ASIC Architecture -- Microarchitecture design, power estimation, process node analysis, migration strategy from FPGA.

    \item[Chapter 12:] Conclusion and Future Work -- Summary of contributions, limitations, and directions for future research.
\end{description}